\documentclass[12pt, a4paper, oneside]{article}

\usepackage[utf8]{inputenc}
\usepackage[T2A]{fontenc}
\usepackage[russian]{babel}

\usepackage[left=1.5cm, right=1.5cm, top=2cm, bottom=2cm]{geometry}

\usepackage{amssymb}
\usepackage{amsmath}
\usepackage{graphicx}
\usepackage{listings}
\usepackage{float}
\usepackage{url}
\usepackage{indentfirst}
\usepackage{setspace}
\usepackage{booktabs}
\usepackage{multirow}
\usepackage{float}

\onehalfspacing

\usepackage{caption}
\captionsetup[figure]{font=small, labelfont=bf}
\captionsetup[table]{font=small, labelfont=bf}

\usepackage[backend=biber,style=numeric]{biblatex}
\addbibresource{bibliography.bib}

\renewcommand{\baselinestretch}{1.4}

\newcommand{\MyUniversity}{Национальный исследовательский университет \\ "Высшая школа экономики"}
\newcommand{\MyProjectText}{Исследовательский проект}
\newcommand{\MyTopicText}{Обучение ретривера на инструкциях}
\newcommand{\MyTopicTextEn}{Promptriever: Instruction-Trained Retrievers}
\newcommand{\MyStudentGroup}{БПМИ242}
\newcommand{\MyStudentName}{Хасанов Айдар Рифкатович}
\newcommand{\MySupervisorName}{Андреева Дарья Александровна}
\newcommand{\MyCityYear}{Москва, 2026}

\newcommand{\maketitlepage}{
    \thispagestyle{empty}
    \begin{center}
        \vspace*{\fill}
        
        \vspace{0.5em}
        {\Large \MyUniversity\\}
        \vspace{4cm}
        {\Large \textbf{\MyProjectText}\\}
        {\large \MyTopicText\\}
        {\large \MyTopicTextEn\\}
        \vspace{2cm}
        
        \begin{minipage}{0.9\textwidth}
            \flushright
            \large
            \textbf{Выполнил:}\\
            Студент группы \MyStudentGroup\\
            \MyStudentName
        \end{minipage}
        
        \vspace{3em}

        \begin{minipage}{0.9\textwidth}
            \flushright 
            \large
            \textbf{Научный руководитель:}\\
            \MySupervisorName
        \end{minipage}
        
        \vspace{4cm}
        {\large \MyCityYear}
        \vspace*{\fill}
    \end{center}
    \newpage
}

% promptbox
\usepackage[most]{tcolorbox}
\usepackage{xcolor}

\definecolor{PromptBg}{HTML}{F4F8FC}
\definecolor{PromptFrame}{HTML}{7A9EBE}
\definecolor{PromptTitleBg}{HTML}{DCEAF7}
\definecolor{PromptTitleText}{HTML}{2F4F6F}

\newtcolorbox{promptnegbox}{
  enhanced,
  breakable,
  colback=PromptBg,
  colframe=PromptFrame,
  boxrule=1pt,
  arc=1.5mm,
  left=4mm,
  right=3mm,
  top=3mm,
  bottom=3mm,
  title={Промпт для генерации негативных инструкций},
  coltitle=PromptTitleText,
  fonttitle=\bfseries,
  colbacktitle=PromptTitleBg,
  attach boxed title to top left={xshift=3mm,yshift=-2mm},
  boxed title style={
    boxrule=0.6pt,
    colframe=PromptFrame,
    arc=1mm
  }
}

\newtcolorbox{promptpassage}{
  enhanced,
  breakable,
  colback=PromptBg,
  colframe=PromptFrame,
  boxrule=1pt,
  arc=1.5mm,
  left=4mm,
  right=3mm,
  top=3mm,
  bottom=3mm,
  title={Промпт для генерации инструктивных негативов},
  coltitle=PromptTitleText,
  fonttitle=\bfseries,
  colbacktitle=PromptTitleBg,
  attach boxed title to top left={xshift=3mm,yshift=-2mm},
  boxed title style={
    boxrule=0.6pt,
    colframe=PromptFrame,
    arc=1mm
  }
}

\definecolor{PromptPosBg}{HTML}{F3FAF4}
\definecolor{PromptPosFrame}{HTML}{7FA98A}
\definecolor{PromptPosTitleBg}{HTML}{DDF0E1}
\definecolor{PromptPosTitleText}{HTML}{355E3B}

\newtcolorbox{promptposbox}{
  enhanced,
  breakable,
  colback=PromptPosBg,
  colframe=PromptPosFrame,
  boxrule=1pt,
  arc=1.5mm,
  left=4mm,
  right=3mm,
  top=3mm,
  bottom=3mm,
  title={Промпт для генерации позитивных инструкций},
  coltitle=PromptPosTitleText,
  fonttitle=\bfseries,
  colbacktitle=PromptPosTitleBg,
  attach boxed title to top left={xshift=3mm,yshift=-2mm},
  boxed title style={
    boxrule=0.6pt,
    colframe=PromptPosFrame,
    arc=1mm
  }
}

\newtcolorbox{datasetexample}{
  enhanced,
  breakable,
  colback=PromptBg,
  colframe=PromptFrame,
  boxrule=1pt,
  arc=1.5mm,
  left=4mm,
  right=3mm,
  top=3mm,
  bottom=3mm,
  title={Пример из обучающего датасета},
  coltitle=PromptTitleText,
  fonttitle=\bfseries,
  colbacktitle=PromptTitleBg,
  attach boxed title to top left={xshift=3mm,yshift=-2mm},
  boxed title style={
    boxrule=0.6pt,
    colframe=PromptFrame,
    arc=1mm
  }
}
% promptbox


\begin{document}

\maketitlepage

\pagenumbering{roman}

\tableofcontents
\newpage

\section*{Аннотация}

\addcontentsline{toc}{section}{Аннотация} 

Данная работа посвящена реализации подхода Promptriever для русского языка путем дообучения нескольких моделей на модифицированном датасете SberQuAD. Целью исследования является создание поисковой системы, способной эффективно следовать текстовым инструкциям пользователя для повышения точности и управляемости поиска. Эффективность метода оценивается на бенчмарках ruMTEB и mFollowIR с использованием метрик nDCG@10 и pMRR.

\newpage

\pagenumbering{arabic}
\setcounter{page}{1}

\section{Введение}
\label{sec:vvedenie}
Поиск информации прошел путь от простого сопоставления ключевых слов до использования семантических векторов, создаваемых нейронными сетями. Современные би-энкодеры позволяют находить документы по смыслу, переводя запросы и документы в единое векторное пространство $\mathbb{R}^d$. Однако традиционные модели поиска ограничены: они не способны учитывать специфические инструкции пользователя, воспринимая их как часть поискового запроса.

Данная работа посвящена реализации подхода, использованного в статье Promptriever \cite{weller2024promptriever} для русского языка. Основная идея заключается в обучении модели воспринимать текстовые инструкции. 

\subsection{Актуальность}

В последние годы методы информационного поиска претерпели значительную трансформацию благодаря внедрению плотных векторных представлений 
на базе предобученных языковых моделей. Несмотря на высокую эффективность в классических задачах поиска по сходству, 
современные ретриверы обладают существенным ограничением: они воспринимают поисковый запрос как единый семантический вектор, 
игнорируя специфические инструкции пользователя.

В реальных сценариях запросы часто содержат не только описание темы, но и дополнительные ограничения. 
Обычные би-энкодеры, такие как multilingual-e5, зачастую не способны 
отличить релевантный документ от документа, который близок по теме, но нарушает заданную инструкцию. 
Это явление ограничивает гибкость поисковых систем, делая их нечувствительными к контексту задачи.

Актуальность данной работы обусловлена необходимостью разработки и имплементации подходов, которые 
позволяют адаптировать ретриверы к следованию инструкциям. 
Это критически важно для создания интеллектуальных ассистентов и систем RAG (Retrieval-Augmented Generation), 
где точность следования промпту напрямую определяет качество итогового ответа системы.

\subsection{Цели}

\begin{enumerate}
    \item Изучение и реализация модели Promptriever, которая умеет работать с инструкциями на естественном языке при выполнении поиска документов. 
    
    \item Оценка эффективности модели в различных условиях. 

    \item Сравнение различных моделей в качестве baseline при реализации ретривера. 
\end{enumerate}

\subsection{Задачи проекта}

\begin{enumerate}
    \item Изучение архитектуры, предложенной в статье Promptriever.
    
    \item Подготовка обучающего датасета на основе SberQuAD \cite{sberquad}  с генерацией позитивных и негативных инструкций для обучения модели следованию командам.
    
    \item Дообучение multilingual-e5-base \cite{wang2024multilingual} и других би-энкодеров. 
    
    \item Оценка качества базового поиска на бенчмарках ruMTEB \cite{snegirev2025russianfocusedembeddersexplorationrumteb} (RuBQ, RiaNews, MIRACL) по метрике nDCG@10.
    
    \item Оценка способности модели следовать инструкциям на retrieval-бенчмарке mFollowIR \cite{weller2025mfollowirmultilingualbenchmarkinstruction} с использованием метрики pMRR.

    \item Сравнение разных моделей в качестве baseline.
\end{enumerate}

\newpage

\section{Обзор и анализ источников}
\label{sec:links}

В основу исследования легли следующие работы:

\subsection{Promptriever}

Статья Promptriever: Instruction-Trained Retrievers Can Be Prompted Like Language Models \cite{weller2024promptriever} является основной 
теоретической и методологической опорой данной работы. В ней вводится подход к обучению ретриверов, при котором модель должна учитывать не 
только семантическую близость запроса и документа, но и дополнительные условия, заданные в текстовой инструкции. Это отличает Promptriever от 
классических би-энкодеров: обычная модель поиска стремится приблизить вектор запроса к вектору тематически релевантного документа, тогда как Promptriever 
дополнительно учится различать документы, которые подходят по теме, но нарушают ограничение из промпта.

Promptriever обучается на инструкциях разного вида, поэтому может использоваться в zero-shot сценариях, 
где заранее неизвестны ни формулировка пользовательского запроса, ни тип ограничения. Особую роль в обучении играют 
негативные примеры: авторы используют документы, которые семантически близки к запросу, но не удовлетворяют инструкции. Такая постановка 
заставляет модель обращать внимание на значимые слова в условии и приводит к заметному росту метрики pMRR: в статье сообщается о приросте на 
$+14.3$ по сравнению с RepLLaMA. Кроме того, авторы показывают масштабируемость подхода: качество следования инструкциям повышается при увеличении 
размера базовой модели и объема синтетических данных.

\subsection{InstructOR}

Работа One Embedder, Any Task: Instruction-Finetuned Text Embeddings \cite{su2023embeddertaskinstructionfinetunedtext} также 
рассматривает идею использования инструкций при построении эмбеддингов. InstructOR обучается переводить в векторное пространство 
не только исходный текст, но и описание задачи, что делает одну и ту же модель применимой к разным сценариям. 
В этом смысле статья важна для данного исследования как более ранний пример instruction-tuning для embedding-моделей.

Однако между InstructOR и Promptriever есть существенное различие. InstructOR обучался на относительно 
фиксированных инструкциях вида \texttt{represent the \{domain\} for \{task type\} supporting documents; input: \{input\}}, 
тогда как Promptriever рассчитан на более свободные и сложные формулировки. Кроме того, Promptriever использует при дообучении 
негативные инструкции к запросам, благодаря чему модель учится отличать документы, которые релевантны теме, от документов, которые 
одновременно релевантны теме и удовлетворяют условию. Это отражается и в качестве: в статье о Promptriever средний pMRR на FollowIR 
составляет $+11.2$ против $-2.8$ у Instructor XL, а на MS-MARCO значение nDCG у Promptriever также существенно выше: 
$92.1$ против $48.6$ \cite{weller2024promptriever}.

При этом подход InstructOR сохраняет свои преимущества. Он хорошо обобщается на 
широкий набор задач, включая классификацию, retrieval и оценку генерации текста. Также инструкции 
в InstructOR были собраны вручную, а не сгенерированы синтетически, что может повышать их качество и 
снижать зависимость результата от ошибок LLM. 

\subsection{Multilingual-E5}

Технический отчет Multilingual E5 Text Embeddings \cite{wang2024multilingual} 
описывает модель multilingual-e5-base, которая используется в данной работе в 
качестве одной из базовых моделей для дообучения. E5 относится к семейству плотных embedding-моделей и 
предназначена для сопоставления текстов в общем векторном пространстве. Для настоящего исследования 
она особенно важна, поскольку изначально обучалась как мультиязычная модель и 
поддерживает 93 языка, включая русский.

Обучение multilingual-e5-base состоит из нескольких этапов. Сначала модель проходит 
контрастивное предобучение на $10^9$ парах вида текст-текст из открытых источников, 
а затем дообучается на $1.6 \cdot 10^6$ текстах и $5 \cdot 10^5$ инструкциях, 
сгенерированных GPT-3.5/4. При этом E5 использует специальный формат входных данных: запросы и 
документы подаются как \texttt{query: \{query\}} и \texttt{passage: \{passage\}}. Это необходимо 
учитывать при подготовке обучающего датасета и при последующем тестировании, поскольку несоблюдение формата 
может ухудшить качество эмбеддингов. Размер модели составляет $278$M параметров, что делает ее 
достаточно компактной для экспериментов, но при этом сильной базовой точкой для русскоязычного поиска.

\subsection{BGE-M3}

Статья M3-Embedding \cite{chen2025m3embeddingmultilingualitymultifunctionalitymultigranularity} представляет 
BGE-M3 --- многофункциональную мультиязычную embedding-модель нового поколения. 
Она поддерживает более 100 языков, включая русский, поэтому подходит для сравнения с 
multilingual-e5-base в рамках данной работы. В отличие от моделей, ориентированных только на 
dense retrieval, BGE-M3 объединяет несколько режимов поиска: dense retrieval, sparse retrieval и 
multi-vector retrieval. Это делает модель более гибкой, так как она может использовать как плотные 
векторные представления, так и более детальные механизмы сопоставления текста.

Для данного исследования BGE-M3 интересна также с 
практической стороны. В отличие от multilingual-e5-base, она не требует 
обязательного форматирования входа в виде \texttt{query: \{query\}}, что упрощает 
подготовку данных и снижает риск ошибки при инференсе. Кроме того, модель поддерживает 
контекст до 8192 токенов, поэтому может применяться для поиска по длинным документам без 
предварительного разбиения на короткие фрагменты. При этом BGE-M3 крупнее multilingual-e5-base: 
ее размер составляет $569$M параметров. В работе она используется как сильный baseline, с 
которым можно сравнивать качество следования инструкциям и качество обычного поиска.

\subsection{SberQuAD}

Работа SberQuAD: Russian Reading Comprehension Dataset: 
Description and Analysis \cite{sberquad} посвящена 
созданию и анализу первого крупного русскоязычного датасета 
для задачи смыслового чтения. SberQuAD является русскоязычным аналогом 
SQuAD и содержит более $5 \cdot 10^4$ троек вида ``контекст, вопрос, ответ'', 
собранных с помощью краудсорсинга на базе статей Википедии.

В данном исследовании SberQuAD используется не как датасет 
для question answering в классической постановке, а как основа для построения 
обучающих пар для ретривера. Наличие контекста и связанного с ним вопроса позволяет 
естественным образом сформировать задачу поиска релевантного фрагмента по запросу. Далее 
такие пары можно модифицировать с помощью позитивных и негативных инструкций, превращая исходный 
QA-датасет в материал для обучения модели следованию пользовательским ограничениям.

\subsection{mFollowIR}

Статья mFollowIR \cite{weller2025mfollowirmultilingualbenchmarkinstruction} вводит 
мультиязычный бенчмарк для оценки того, насколько хорошо ретриверы следуют инструкциям. 
Для данной работы этот источник особенно важен, потому что mFollowIR включает русский язык 
и позволяет измерять не только тематическую релевантность выдачи, но и выполнение 
дополнительных условий, заданных пользователем.

Бенчмарк содержит запросы с ограничениями разного типа: например, 
по временному интервалу, стилю текста или подмножеству сущностей. 
Такая постановка ближе к реальным поисковым сценариям, где пользователь 
часто ищет не просто документы по теме, а документы, удовлетворяющие конкретному условию. 
Для оценки используется метрика pMRR, которая штрафует модель, если в верхние позиции выдачи 
попадает документ, релевантный теме, но нарушающий инструкцию. Поэтому ситуация, в которой 
модель имеет высокий nDCG и низкий pMRR, означает, что она умеет находить тематически 
близкие документы, но плохо следует инструкциям.

\subsection{ruMTEB}

Работа The Russian-focused embedders' exploration: ruMTEB benchmark 
and Russian embedding model design \cite{snegirev2025russianfocusedembeddersexplorationrumteb} 
представляет ruMTEB --- русскоязычное расширение глобального бенчмарка MTEB. 
Авторы объединяют 23 датасета, разделенных на 7 категорий задач, включая классификацию, 
кластеризацию и поиск. Благодаря этому ruMTEB позволяет оценивать embedding-модели не на 
одной узкой задаче, а на широком наборе сценариев, характерных для русского языка.

В рамках данной работы ruMTEB используется прежде всего для оценки базового 
качества поиска. Для retrieval-задач бенчмарк включает такие наборы данных, 
как RuBQ, RiaNews и русскоязычная версия MIRACL. 
Эти датасеты позволяют проверить, сохраняет ли модель качество обычного семантического 
поиска после дообучения на инструкциях. Такое разделение оценивания важно: mFollowIR 
показывает способность модели следовать условиям, а ruMTEB помогает убедиться, 
что улучшение instruction following не достигается ценой деградации 
стандартного retrieval-качества.

\newpage

\section{Основная часть}
\label{sec:main}

\subsection{Первые эксперименты}

На начальном этапе работы была проведена оценка моделей multilingual-e5-base и bge-m3 без дообучения на инструкциях. Тестирование проводилось на датасете RuBQRetrieval (часть бенчмарка ruMTEB) и mFollowIR.

\begin{table}[H]
    \centering
    \caption{Результаты тестирования без дообучения}
    \label{tab:instruction_results}
    \begin{tabular}{lccc}
        \toprule
        \textbf{Модель} & \textbf{RuBQRetrieval (nDCG@10)} & \textbf{mFollowIR (pMRR)} \\
        \midrule
        multilingual-e5-base\textsuperscript{1} & $69.170$ & $-2.994$ \\
        bge-m3 & $71.178$ & $-2.154$ \\
        \bottomrule
    \end{tabular}
    \vspace{0.2cm}
\end{table}

\footnotetext[1]{При тестировании использовался формат \texttt{query: \{запрос\}} и \texttt{passage: \{текст\}}}

Как видно из таблицы, без дообучения модели показывают неплохие результаты по метрике nDCG, но достаточно слабые результаты по метрике pMRR. Дальнейшие эксперименты будут направлены на обучение модели таким образом, чтобы при сохранении высокого качества поиска она демонстрировала значительный прирост по метрике pMRR. 

\subsection{Сбор тренировочного датасета}

За основу датасета для дообучения на инструктивных негативах берется 
SberQuAD. Для каждого примера (контекст, вопрос, ответ) генерируются позитивные и негативные инструкции. 
Позитивные инструкции формулируются так, чтобы модель должна была выдавать релевантные документы, 
а негативные инструкции содержат условия, которые нарушают релевантность. 

Инструкции генерировались синтетически с помощью LLM. В 
качестве генератора использовалась llama3-8b \cite{llama3modelcard}.

Разберем генерацию датасета по шагам. 
По паре $(\texttt{query, context})$, где
$\texttt{context}$ -- отрывок из текста (взятый из исходного датасета SberQuAD), 
$\texttt{query}$ -- вопрос (также взятый из SberQuAD), строилась пара $(\texttt{query + instruction}, 
\texttt{context})$, $\texttt{instruction}$ 
-- сгенерированная инструкция. 

\begin{promptposbox}
\small
Ты помогаешь строить обучающий датасет для instruction-following retrieval.
По паре вопрос-контекст нужно придумать ОДНУ позитивную инструкцию на русском языке.
Позитивная инструкция должна:
\begin{enumerate}
    \item звучать как реальная поисковая инструкция;
    \item быть тематически близкой к вопросу;
    \item делать данный контекст ПОДХОДЯЩИМ и релевантным;
    \item не должна зависеть от скрытой разметки;
    \item не должна быть бессмысленной, слишком общей или противоречивой.
\end{enumerate}
Верни JSON с полями: \texttt{positive\_instruction}, \texttt{relevance\_reason}.
\end{promptposbox}

Затем генерировался текст, который делал $\texttt{context}$ 
релевантным для $\texttt{query + instruction}$ (в статье такой текст называют instruction positive). 
И таким же образом 3 текста, которые подходят под $\texttt{query}$, но не подходят под 
$\texttt{instruction}$ (в статье такие называют instruction negative). 

\begin{promptpassage}
\small
Ты помогаешь строить instruction-following retrieval датасет по рецепту Promptriever.
Для каждого примера нужно создать passages, которые позволяют обучать bi-encoder
учитывать инструкцию на уровне конкретного запроса.

Верни JSON со следующими полями:
\begin{enumerate}
    \item \texttt{generated\_positive\_passage}: passage, релевантный и вопросу, и инструкции
    \item \texttt{instruction\_negative\_passages}: список из нескольких passages, которые релевантны исходному вопросу, 
    но НЕ удовлетворяют дополнительному условию инструкции
\end{enumerate}
Требования:
\begin{enumerate}
    \item passages должны быть правдоподобными и естественными;
    \item нельзя делать negative явно абсурдным или нерелевантным теме;
    \item negative должен отвечать базовому вопросу, но не проходить по instruction;
    \item верни только JSON без markdown.
\end{enumerate}
\end{promptpassage}

После этого к каждому $\texttt{query}$ искали 2 текста
из исходного датасета, которые близки по косинусному расстоянию, но не подходят 
по смыслу (в статье такие называют hard negative). 
Все эти тексты добавляются в датасет. 

\subsubsection*{Валидация сгенерированных текстов}
После генерации текстов, их нужно валидировать, потому что 
LLM могла сгенерировать текст, который не будет 
релевантен запросу с инструкцией, хотя должен быть (или наоборот). 
Для валидации мы используем bge-reranker-m3-v2 \cite{chen2025m3embeddingmultilingualitymultifunctionalitymultigranularity}.
Каждому тексту и запросу (возможно, с инструкцией) 
сопоставляет скор -- меру сходства. Инструкции и сгенерированные запросы 
с низким скором отбрасываются. 

\begin{table}[H]
    \centering
    \caption{Пороги для сгенерированных текстов / инструкций}
    \label{tab:generation_thresholds}
    \begin{tabular}{cccc}
        \toprule
        \textbf{Инструкция} & \textbf{Instruction positive текст} & \textbf{Instruction negative текст} \\
        \midrule
        $\ge 1.0$ & $\ge 3.0$ & $\le 4.0$ \\
        \bottomrule
    \end{tabular}
    \vspace{0.2cm}
\end{table}

К первым $5745$ строкам из SberQuAD были сгенерированы позитивные инструкции, 
$5062$ из них были помечены как релевантные. 
После валидации instruction positive текстов (для каждого запроса генерировалось по $3$ текста) 
$5219$ были отмечены как релевантные. После валидации instruction negative $4193$ текста были отмечены как релевантные. 

Итоговый размер датасета -- $1748$ строк. В каждой из них было $\ge 1$ hard negatives, $\le 3$ instruction negatives и 
$1$ instruction positive текст. 

\begin{datasetexample}
    \textbf{Query:} Шедеврами какой музыки признаны сочинения Моцарта? \\
    \textbf{Instruction:} Найдите информацию о музыкальных жанрах, в которых признаны шедевры Моцарта.
    
    \tcbline
    
    \textbf{Positive:} Моцарт - один из наиболее значительных композиторов Венской классической школы, и его сочинения признаны шедеврами в различных музыкальных жанрах, включая симфоническую, концертную, камерную, оперную и хоровую музыку.
    
    \medskip
    \textbf{Instruction Negative:} В музыке Моцарта можно услышать влияние классицизма и барокко...
    
    \medskip
    \textbf{Hard Negative:} Отличительной чертой творчества Моцарта является сочетание строгих, ясных форм...
\end{datasetexample}

\subsection{Дообучение модели}

Дообучение модели производилось с помощью LoRA. 
Было проведено несколько экспериментов с разными моделями и разными параметрами LoRA. 
В качестве loss-функции использовалась MultipleNegativesRankingLoss. 
Каждая модель дообучалась ровно $20$ эпох. 

\begin{table}[H]
    \centering
    \caption{Результаты дообучения с различными параметрами LoRA}
    \label{tab:lora_results}
    \begin{tabular}{lccccccc}
        \toprule
        \textbf{Base model} & $r$ & $\alpha$ & \textbf{Dropout} & \textbf{Batch size} & \textbf{Trainable params\textsuperscript{1}} & \textbf{nDCG@10} & \textbf{pMRR} \\
        \midrule
        m-e5-base & 8 & 16 & 0 & 16 & 0.16\% (0.45M) & $59.405$ & $-0.08538$ \\
        m-e5-base & 8 & 16 & 0.05 & 16 & 0.16\% (0.45M) & $61.392$ & $-0.09782$ \\
        m-e5-base & 16 & 32 & 0.05 & 16 & 0.32\% (0.89M) & $59.894$ & $0.61896$ \\
        m-e5-base & 32 & 64 & 0.05 & 16 & 0.63\% (1.75M) & $59.785$ & $-0.07843$ \\
        \midrule
        bge-m3 & 8 & 16 & 0.05 & 16 & 0.21\% (1.12M) & $63.579$ & $-0.08945$ \\
        bge-m3 & 16 & 32 & 0.05 & 16 & 0.41\% (2.23M) & $61.962$ & $-0.05129$ \\
        \bottomrule
    \end{tabular}
    \vspace{0.2cm}
\end{table}

\footnotetext[1]{Доля от всех параметров модели}

После экспериментов с дообучением получили прирост в $>+3$ по метрике 
pMRR относительно baseline решения, однако также получили 
заметную просадку в $-10$ по метрике nDCG@10 (в среднем). 
Это подтверждает гипотезу, что модель научилась различать 
некоторые дополнительные инструкции в запросах и менять 
эмбеддинг в зависимости от них, однако при этом ее 
функции ретрива немного ухудшились. 
На это могло повлиять несколько факторов, как то:

\begin{enumerate}
    \item размер и качество синтетического датасета; есть риск, что мы плохо отсортировали сгенерированные тексты, 
    и в обучающий датасет попало много шумных примеров, что могло ухудшить качество обучения.
    
    \item параметры LoRA; возможно следовало провести более тщательный подбор гиперпараметров, 
    например, попробовать другой learning rate, большее число эпох, другие значения $r$ и так далее. 

    \item архитектура модели; возможно, базовая модель не была достаточно мощной 
    для эффективного обучения на инструкциях, и более крупная 
    модель могла бы показать лучший результат.
\end{enumerate}

Лучшая попытка показала прирост в $+3.6$ по метрике pMRR, 
что является неплохим результатом, учитывая ограниченный размер датасета и 
относительно слабую модель, использовавшуюся при его генерации.

\subsection{Latency}

Помимо качества поиска, для retrieval-модели важно время построения эмбеддинга. 
Поэтому была добавлена отдельная проверка latency: для одного запроса выполнялось несколько прогонов 
после warmup, после чего считались среднее значение, медиана и p95 ($95\%$-й перцентиль).

\begin{table}[H]
    \centering
    \caption{Latency кодирования одного запроса. Замер проводился на RTX4090.}
    \label{tab:latency_results}
    \begin{tabular}{lccc}
        \toprule
        \textbf{Модель}& \textbf{Mean, ms} & \textbf{Median, ms} & \textbf{p95, ms} \\
        \midrule
        multilingual-e5-base & $6.913$ & $6.798$ & $7.409$ \\
        m-e5-base + LoRA ($r=16, \alpha=32$) & $6.861$ & $6.805$ & $7.441$ \\
        bge-m3 & $12.659$ & $12.552$ & $13.793$ \\
        bge-m3 + LoRA ($r=16, \alpha=32$) & $12.448$ & $12.262$ & $14.046$ \\
        \bottomrule
    \end{tabular}
    \vspace{0.2cm}
\end{table}

Замер latency проводился следующим образом: для каждой модели выполнялось $200$ прогонов 
кодирования одного запроса с инструкцией, первые $20$ прогонов -- warmup, а затем для оставшихся 
$180$ прогонов считались среднее значение, медиана и p95.

Можем видеть, что время инференса практически не изменилось между 
дообученной моделью и базовой. 
Если сравнивать между собой модели на основе bge-m3 и m-e5-base, то 
bge-m3 работает примерно в 2 раза медленнее, 
что связано с ее большим размером ($569$M параметров против $278$M). 

\newpage

\section{Заключение}

В рамках работы был реализован пайплайн дообучения retrieval-модели на русскоязычных 
инструкциях по мотивам подхода Promptriever. Был собран датасет на основе SberQuAD: 
для исходных вопросов генерировались позитивные инструкции, instruction positive тексты, 
instruction negative тексты и hard negatives. Дополнительно была проведена автоматическая 
валидация сгенерированных примеров с помощью reranker-модели, чтобы уменьшить количество 
шумных обучающих пар.

Эксперименты показали, что дообучение с помощью LoRA может заметно улучшить качество 
следования инструкциям. Наиболее удачный запуск для multilingual-e5-base дал прирост 
по метрике pMRR относительно baseline при сравнительно небольшой просадке в nDCG@10. 
Это означает, что модель стала лучше отличать документы, которые просто близки к запросу 
по теме, от документов, которые действительно удовлетворяют дополнительному условию 
из инструкции.

При этом результаты остаются ограниченными размером и качеством синтетического датасета. 
Итоговый набор данных получился сравнительно небольшим, а генерация инструкций и негативов 
не всегда идеально отражает реальные пользовательские сценарии. Поэтому полученные результаты 
следует рассматривать как подтверждение применимости подхода.

\newpage

\section{Использование ИИ}

Во время работы над проектом использовались ИИ-инструменты 
для написания, отладки и оптимизации кода. Для этих целей использовался ChatGPT. 
Все остальное -- поиск и анализ литературы, проектирование экспериментов, интерпретация результатов -- 
выполнялось самостоятельно.

\newpage

\section{Дальнейшая работа}

В дальнейшем имеет смысл увеличить размер обучающего датасета и провести более строгую 
фильтрацию синтетических примеров. Сейчас обучение проводится на ограниченной части SberQuAD, 
поэтому модель видит только часть возможных типов инструкций. Расширение датасета позволит 
проверить, растет ли качество instruction following при увеличении числа примеров.

Также важно провести более подробные ablation-эксперименты. В частности, отдельно сравнить 
вклад hard negatives и instruction negatives, проверить разные значения ранга LoRA, число эпох 
и learning rate. Это позволит понять, какие компоненты пайплайна действительно дают прирост, 
а какие почти не влияют на итоговое качество.

Еще одно направление -- сравнение нескольких backbone-моделей. В данной работе основное внимание 
уделялось multilingual-e5-base, однако bge-m3 потенциально лучше подходит для русского языка 
и multilingual retrieval. Поэтому отдельный полноценный запуск bge-m3 с тем же набором метрик 
позволит понять, насколько результат зависит от выбранной базовой модели.

Наконец, полезно добавить ручную оценку части сгенерированных инструкций и негативов. 
Автоматическая валидация через reranker удобна для масштабирования, но она не всегда замечает 
логические ошибки в инструкции. Ручная проверка небольшой выборки помогла бы точнее оценить 
качество датасета и объяснить ошибки модели на mFollowIR.

\newpage

\printbibliography[heading=bibintoc, title={Список использованных источников}]

\end{document}
